\documentclass[11pt]{article}

\usepackage[margin=1in]{geometry}
\usepackage[T1]{fontenc}
\usepackage[utf8]{inputenc}
\usepackage{lmodern}
\usepackage{microtype}
\usepackage{amsmath,amssymb,amsthm,mathtools,bm}
\usepackage{enumitem}
\usepackage{booktabs}
\usepackage{longtable}
\usepackage{xcolor}
\usepackage{hyperref}
\usepackage[nameinlink,noabbrev]{cleveref}

\hypersetup{
  colorlinks=true,
  linkcolor=blue!60!black,
  citecolor=blue!60!black,
  urlcolor=blue!60!black,
  pdftitle={Minimum Useful Simulation for a Control-Anchored Mean-Field Neural SDE Model},
  pdfauthor={OpenAI}
}

\title{Minimum Useful Simulation for a Control-Anchored\\Mean-Field Neural SDE Model of $P4 \to P60$ Perturb-seq Dynamics}
\author{}
\date{}

\newcommand{\R}{\mathbb{R}}
\newcommand{\E}{\mathbb{E}}
\newcommand{\Var}{\mathrm{Var}}
\newcommand{\Prob}{\mathbb{P}}
\newcommand{\dd}{\,\mathrm{d}}
\newcommand{\KL}{\mathrm{KL}}
\newcommand{\diag}{\mathrm{diag}}
\newcommand{\tr}{\mathrm{tr}}
\newcommand{\simplex}{\Delta}
\newcommand{\Gcal}{\mathcal{G}}
\newcommand{\Gctrl}{\mathcal{G}_{\mathrm{ctrl}}}
\newcommand{\Scal}{\mathcal{S}}
\newcommand{\Zcal}{\mathcal{Z}}
\newcommand{\Mplus}{\mathcal{M}_{+}}

\newcommand{\norm}[1]{\left\lVert #1 \right\rVert}
\newcommand{\abs}[1]{\left\lvert #1 \right\rvert}
\newcommand{\inner}[2]{\left\langle #1, #2 \right\rangle}

\begin{document}
\maketitle

\begin{abstract}
This document gives an explicit introduction to the \emph{minimum useful simulation} for a control-anchored mean-field neural stochastic differential equation model of perturbation-conditioned single-cell dynamics over the interval $P4 \to P60$. The central idea is that a simulation benchmark should be only as complicated as necessary to test the scientific claims of the model. A benchmark that is too simple fails to test key components such as separate drift, diffusion, and growth effects. A benchmark that is too complicated makes debugging impossible and confounds modeling failure with simulator complexity. For this reason, the minimum useful simulation should be built in two stages. The first stage is a core one-dimensional benchmark without mean-field feedback, designed to test control anchoring, perturbation-specific drift, perturbation-specific diffusion, perturbation-specific reaction or growth, and finite-measure endpoint fitting. The second stage is a minimal mean-field benchmark with two synthetic screens and one context-driving perturbation, designed to test whether the model truly benefits from a population-feedback layer rather than merely containing one in notation. The document gives the full simulator specification, explicit parameter values, sampling scheme, expected outcomes, and ablation logic.
\end{abstract}

\tableofcontents

\section{Why a minimum useful simulation is needed}

A simulation study for this modeling framework should not aim to reproduce every biological detail of an in vivo perturb-seq experiment. Its purpose is more basic: it should test whether the mathematical architecture of the proposed model is necessary, identifiable enough to be learned under endpoint-only supervision, and internally consistent.

A useful simulation therefore has to satisfy two competing requirements:
\begin{enumerate}[label=(\roman*),leftmargin=2.2em]
    \item it must be \emph{simple enough} that one can understand what each model component is doing and diagnose failure unambiguously;
    \item it must be \emph{rich enough} that each major layer of the model has a genuine reason to exist.
\end{enumerate}

For the present formulation, the main architectural claims are:
\begin{enumerate}[label=(C\arabic*),leftmargin=2.6em]
    \item perturbations should be modeled relative to a shared \emph{control anchor};
    \item state evolution should contain both \emph{drift} and \emph{diffusion};
    \item abundance change should be represented by a separate \emph{reaction or growth field};
    \item endpoint supervision should act on a \emph{finite measure}, not only on a normalized state distribution;
    \item a \emph{mean-field context layer} should be used only if it explains variation that cannot be absorbed by a purely perturbation-conditioned but context-free dynamics.
\end{enumerate}

No single tiny benchmark can test all of these equally well. The right design is therefore a two-stage ladder:
\begin{enumerate}[label=(S\arabic*),leftmargin=2.6em]
    \item a \textbf{core benchmark} that tests the control-anchored weighted SDE without mean-field feedback;
    \item a \textbf{mean-field benchmark} that minimally extends the core benchmark in a way that makes shared context genuinely necessary.
\end{enumerate}

\section{Design principles}

The minimum useful simulation should obey the following principles.

\subsection{Use a known latent space}

The simulation should take place directly in a synthetic latent space
\[
\Zcal \subseteq \R^d
\]
rather than simulating raw gene-expression counts and then fitting an encoder. Otherwise, failures of the dynamical model become confounded with representation-learning failures.

\subsection{Make the simulator simpler than the fitted model}

The ground-truth dynamics should be linear or nearly linear, even if the fitted model is neural. A good simulator is one whose behavior is analytically interpretable. The fitted model may then be more expressive than the truth, which makes the benchmark conservative.

\subsection{Keep endpoint-only supervision}

The simulator may generate dense trajectories internally, but the fitted model should only see:
\begin{itemize}[leftmargin=2em]
    \item a $P4$ snapshot,
    \item a $P60$ snapshot,
    \item and guide-abundance-scale masses.
\end{itemize}
Otherwise the benchmark no longer matches the intended inference problem.

\subsection{Decouple state-sample size from mass scale}

The numbers of sampled cells should not be used as a proxy for perturbation mass. The benchmark should explicitly maintain separate objects:
\begin{itemize}[leftmargin=2em]
    \item sampled cell states used to define empirical measures;
    \item guide-abundance masses used to define finite-measure scale.
\end{itemize}
This is necessary to test whether finite-measure endpoint fitting is actually doing meaningful work.

\subsection{Add mean-field only when it is identifiable}

A single synthetic screen is usually not enough to test the necessity of a mean-field context layer. In many cases, a single shared context trajectory can be absorbed into a time-varying baseline drift. Therefore, the minimum useful test of mean-field coupling should include at least two synthetic screens whose global context trajectories differ.

\section{Stage I: the minimum useful core benchmark}

\subsection{Purpose}

The purpose of the core benchmark is to test the smallest nontrivial version of the model:
\begin{enumerate}[label=(\roman*),leftmargin=2.2em]
    \item control anchoring;
    \item separate perturbation effects on drift, diffusion, and reaction;
    \item finite-measure endpoint fitting;
    \item recovery of endpoint effects from two destructive snapshots only.
\end{enumerate}

This benchmark deliberately does \emph{not} test mean-field feedback. That omission is intentional. One should first verify that the weighted control-anchored dynamics works without context.

\subsection{State space and perturbations}

Take latent dimension
\[
d=1.
\]
Let the perturbation set be
\[
\Gcal=\{\mathrm{ctrl},\mathrm{drift},\mathrm{diff},\mathrm{react}\},
\qquad
\Gctrl=\{\mathrm{ctrl}\}.
\]
Thus there are four perturbations:
\begin{enumerate}[label=(\arabic*),leftmargin=2.2em]
    \item one control perturbation;
    \item one perturbation that changes \emph{drift} only;
    \item one perturbation that changes \emph{diffusion} only;
    \item one perturbation that changes \emph{reaction} only.
\end{enumerate}

This is the smallest setup in which the three modeled biological mechanisms are separated cleanly.

\subsection{Normalized time and initial condition}

Let normalized time be
\[
t\in[0,1].
\]
For every perturbation \(g\), initialize from the same $P4$ law
\begin{equation}
Z_0^g \sim \mathcal{N}(m_0,v_0),
\qquad
m_0=-1,
\qquad
v_0=(0.15)^2.
\label{eq:init_law}
\end{equation}
Set the initial guide-abundance scale to
\begin{equation}
M_g^0 = 1
\qquad
\text{for all } g\in\Gcal.
\label{eq:initial_mass_scale}
\end{equation}

\subsection{Ground-truth core dynamics}

For each perturbation \(g\), generate the true latent dynamics by the one-dimensional Ornstein--Uhlenbeck reaction system
\begin{equation}
\dd Z_t^g
=
\kappa(\theta_g-Z_t^g)\,\dd t
+
\sigma_g\,\dd W_t,
\label{eq:core_truth_state}
\end{equation}
together with the weight equation
\begin{equation}
\frac{\dd}{\dd t}\log w_t^g
=
\rho_g.
\label{eq:core_truth_weight}
\end{equation}
Equivalently,
\begin{equation}
w_t^g = \exp(\rho_g t).
\label{eq:core_truth_weight_solution}
\end{equation}

Here:
\begin{itemize}[leftmargin=2em]
    \item \(\kappa>0\) is the attraction strength toward a perturbation-specific target \(\theta_g\);
    \item \(\sigma_g>0\) is the perturbation-specific diffusion amplitude;
    \item \(\rho_g\in\R\) is the perturbation-specific reaction or growth rate.
\end{itemize}

\subsection{Explicit parameter values}

Set
\[
\kappa = 2.
\]
Then choose the perturbation-specific parameters as follows:
\begin{align}
\theta_{\mathrm{ctrl}} &= 0,
&
\sigma_{\mathrm{ctrl}} &= 0.15,
&
\rho_{\mathrm{ctrl}} &= 0,
\label{eq:ctrl_params}
\\[0.5ex]
\theta_{\mathrm{drift}} &= 0.6,
&
\sigma_{\mathrm{drift}} &= 0.15,
&
\rho_{\mathrm{drift}} &= 0,
\label{eq:drift_params}
\\[0.5ex]
\theta_{\mathrm{diff}} &= 0,
&
\sigma_{\mathrm{diff}} &= 0.35,
&
\rho_{\mathrm{diff}} &= 0,
\label{eq:diff_params}
\\[0.5ex]
\theta_{\mathrm{react}} &= 0,
&
\sigma_{\mathrm{react}} &= 0.15,
&
\rho_{\mathrm{react}} &= -0.7.
\label{eq:react_params}
\end{align}

These values are chosen so that:
\begin{itemize}[leftmargin=2em]
    \item the drift perturbation changes only the terminal location;
    \item the diffusion perturbation changes only the terminal spread;
    \item the reaction perturbation changes only the terminal mass.
\end{itemize}

\subsection{Why this benchmark is analytically useful}

Because \cref{eq:core_truth_state,eq:core_truth_weight} is linear, its terminal law is known in closed form. Specifically,
\begin{equation}
\E[Z_1^g]
=
\theta_g + (m_0-\theta_g)e^{-\kappa},
\label{eq:terminal_mean}
\end{equation}
\begin{equation}
\Var(Z_1^g)
=
v_0 e^{-2\kappa}
+
\frac{\sigma_g^2}{2\kappa}\big(1-e^{-2\kappa}\big),
\label{eq:terminal_var}
\end{equation}
and the terminal mass is
\begin{equation}
M_g^1
=
M_g^0 e^{\rho_g}
=
e^{\rho_g}.
\label{eq:terminal_mass_truth}
\end{equation}

Thus the simulation has three analytically separable endpoint signatures:
\begin{enumerate}[label=(\roman*),leftmargin=2.2em]
    \item location shift through \(\theta_g\);
    \item spread change through \(\sigma_g\);
    \item mass change through \(\rho_g\).
\end{enumerate}

\subsection{Observed synthetic data}

For each perturbation \(g\), create the observed snapshots as follows.

\paragraph{Step 1: sample the $P4$ snapshot.}
Draw
\[
n_g^0 = 256
\]
independent samples from the initial law \cref{eq:init_law},
\[
x_{g1},\dots,x_{g n_g^0}\sim p_0^g.
\]
This defines the empirical initial state distribution
\begin{equation}
p_0^g
=
\frac{1}{n_g^0}\sum_{i=1}^{n_g^0}\delta_{x_{gi}}.
\label{eq:empirical_p0_core}
\end{equation}

\paragraph{Step 2: sample the $P60$ snapshot.}
Draw
\[
n_g^1 = 256
\]
independent samples from the normalized terminal state law of \cref{eq:core_truth_state},
\[
y_{g1},\dots,y_{g n_g^1}\sim p_1^g.
\]
This defines
\begin{equation}
p_1^g
=
\frac{1}{n_g^1}\sum_{j=1}^{n_g^1}\delta_{y_{gj}}.
\label{eq:empirical_p1_core}
\end{equation}

\paragraph{Step 3: record the terminal guide-abundance mass.}
Use the separately recorded mass
\begin{equation}
M_g^1 = e^{\rho_g}.
\label{eq:observed_mass_core}
\end{equation}

\paragraph{Step 4: define the finite-measure targets.}
Then the synthetic finite-measure endpoint objects are
\begin{equation}
\mu_0^g = M_g^0 p_0^g,
\qquad
\nu_1^g = M_g^1 p_1^g.
\label{eq:finite_measure_core}
\end{equation}

\paragraph{Important point.}
The sampled cell numbers \(n_g^0\) and \(n_g^1\) are fixed across perturbations, while the mass \(M_g^1\) varies by perturbation. This is essential. It forces the model to learn abundance change from the finite-measure target rather than from trivial sample-size differences.

\subsection{What this benchmark should test}

A fitted model passes the core benchmark if it recovers the following qualitative endpoint structure:
\begin{enumerate}[label=(\roman*),leftmargin=2.2em]
    \item the drift perturbation has a larger terminal mean than control;
    \item the diffusion perturbation has a larger terminal variance than control;
    \item the reaction perturbation has smaller terminal mass than control.
\end{enumerate}

Equivalently, the model should learn that different perturbations affect different layers of the dynamics, rather than explaining all endpoint differences by only one mechanism.

\section{Why this is the minimum useful core test}

This benchmark is minimal in a precise sense.

\subsection{Why one latent dimension is enough}

The purpose of the minimum benchmark is not to test geometry learning in a complex manifold. It is to test whether the model can distinguish:
\begin{itemize}[leftmargin=2em]
    \item movement,
    \item spread,
    \item and mass change
\end{itemize}
from endpoint-only supervision. One dimension is sufficient for that.

\subsection{Why four perturbations are enough}

With only a control and one non-control perturbation, one cannot tell whether the model is using drift, diffusion, or reaction correctly. The four-perturbation design is the smallest clean factorial decomposition:
\begin{center}
control \quad / \quad drift-only \quad / \quad diffusion-only \quad / \quad reaction-only.
\end{center}

\subsection{Why fixed cell counts are necessary}

If the number of sampled $P60$ cells were allowed to track \(M_g^1\), then the model could partly infer abundance from the number of observations alone. That would make the finite-measure formulation less meaningful. Fixing the cell counts avoids this shortcut.

\subsection{Why the truth should be linear}

The fitted model may use neural networks, but the simulator should remain interpretable. In the minimum benchmark, the purpose is to test whether the model architecture can recover endpoint effects, not whether it can approximate arbitrary nonlinear dynamics.

\section{Stage II: the minimum useful mean-field benchmark}

\subsection{Why the core benchmark is not enough for mean-field}

The core benchmark deliberately omits mean-field feedback. This is appropriate for debugging the weighted SDE machinery. However, it does not test whether a context layer is useful.

A single synthetic screen with one context trajectory is generally not enough to justify a mean-field model. In many cases, a single shared context path can be absorbed into a flexible time-dependent baseline coefficient. Therefore, a benchmark that truly tests the mean-field claim must generate at least two distinct global-context trajectories.

\subsection{Screens and perturbations}

Introduce a screen index
\[
s \in \Scal = \{1,2\}.
\]
Keep the previous perturbations and add one context-driving perturbation:
\[
\Gcal
=
\{\mathrm{ctrl},\mathrm{drift},\mathrm{diff},\mathrm{react},\mathrm{driver}\}.
\]
The driver perturbation is not introduced to create biological realism in itself. It is introduced to create a perturbation whose abundance changes the global environment experienced by all other perturbations.

\subsection{Latent state and context summary}

Continue to use
\[
d=1.
\]
Define a smooth occupancy function
\begin{equation}
h(z) = \frac{1}{1+e^{-6z}}.
\label{eq:occupancy_function}
\end{equation}
This function is close to \(0\) for states on the left side of the latent axis and close to \(1\) for states on the right side. It therefore serves as a soft right-side occupancy indicator.

For each screen \(s\), define the context summary
\begin{equation}
c_{s,t}
=
\frac{1}{\bar M_s(t)}
\sum_{g\in\Gcal}
\int_{\Zcal} h(z)\,\mu_{s,t}^g(\dd z),
\qquad
\bar M_s(t)=\sum_{g\in\Gcal}\mu_{s,t}^g(\Zcal).
\label{eq:screen_context}
\end{equation}

\subsection{Ground-truth mean-field dynamics}

For each screen \(s\) and perturbation \(g\), generate the latent state by
\begin{equation}
\dd Z_{s,t}^g
=
\Big[
\kappa(\theta_g-Z_{s,t}^g) + \eta c_{s,t}
\Big]\dd t
+
\sigma_g\,\dd W_t,
\label{eq:mf_truth_state}
\end{equation}
and retain the same reaction law
\begin{equation}
\frac{\dd}{\dd t}\log w_{s,t}^g = \rho_g.
\label{eq:mf_truth_weight}
\end{equation}

The only new term relative to the core benchmark is \(\eta c_{s,t}\), which means that the drift experienced by all populations depends on the current screen-level occupancy of the right side of state space.

\subsection{Driver perturbation parameters}

Use the previous parameters for control, drift, diffusion, and reaction. For the driver perturbation, set
\begin{equation}
\theta_{\mathrm{driver}} = 1.0,
\qquad
\sigma_{\mathrm{driver}} = 0.15,
\qquad
\rho_{\mathrm{driver}} = \log 2.
\label{eq:driver_params}
\end{equation}
Set the context-coupling strength to
\begin{equation}
\eta = 0.8.
\label{eq:eta_value}
\end{equation}

Thus the driver tends to move toward the right side of the latent axis and expand in mass, which causes its presence to influence the context experienced by all perturbations.

\subsection{How the two screens differ}

To make the mean-field effect identifiable, the two screens should differ only in their initial abundance of the driver perturbation. Set:
\begin{equation}
M_{\mathrm{driver}}^{0,(1)} = 0.5,
\qquad
M_{\mathrm{driver}}^{0,(2)} = 2.0.
\label{eq:driver_screen_difference}
\end{equation}
For all other perturbations in both screens, set
\begin{equation}
M_g^{0,(1)} = M_g^{0,(2)} = 1.
\label{eq:other_screen_masses}
\end{equation}

This means that the intrinsic perturbation laws are the same in both screens, but the global context trajectory differs because the driver starts at different abundance.

\subsection{Why two screens are enough}

Two screens are sufficient because they induce two different context trajectories:
\begin{itemize}[leftmargin=2em]
    \item screen \(1\) has a weak driver-induced rightward context;
    \item screen \(2\) has a strong driver-induced rightward context.
\end{itemize}
Therefore, even the control perturbation experiences different effective evolution across the two screens despite having the same intrinsic control law.

A model with a sample- or screen-specific context layer can explain this coherently. A context-free model cannot do so without introducing separate nuisance laws for each screen, which defeats the purpose of mean-field sharing.

\subsection{Observed synthetic data for the mean-field benchmark}

For each screen \(s\) and perturbation \(g\):
\begin{enumerate}[label=(\arabic*),leftmargin=2.2em]
    \item sample \(n_{s,g}^0=256\) cells at \(t=0\) from the initial law;
    \item simulate the full weighted system under \cref{eq:mf_truth_state,eq:mf_truth_weight};
    \item sample \(n_{s,g}^1=256\) cells at \(t=1\) from the normalized terminal state law;
    \item record the terminal guide-abundance mass \(M_{s,g}^1\);
    \item define
    \[
    \mu_{0}^{s,g}=M_{s,g}^0p_0^{s,g},
    \qquad
    \nu_1^{s,g}=M_{s,g}^1p_1^{s,g}.
    \]
\end{enumerate}

\section{Training and evaluation protocol}

\subsection{What the fitted model should observe}

In both the core and mean-field benchmarks, the fitted model should observe only:
\begin{itemize}[leftmargin=2em]
    \item the perturbation labels,
    \item the $P4$ cell states,
    \item the $P60$ cell states,
    \item and the guide-abundance masses \(M_g^0,M_g^1\) or \(M_{s,g}^0,M_{s,g}^1\).
\end{itemize}

The model should \emph{not} observe the true intermediate trajectories or the true coefficients \((\theta_g,\sigma_g,\rho_g)\).

\subsection{Primary evaluation quantities}

The following endpoint summaries should be compared between truth and fit:
\begin{enumerate}[label=(E\arabic*),leftmargin=2.6em]
    \item terminal mean:
    \[
    \hat m_g^1 = \int z\,\hat p_1^g(\dd z);
    \]
    \item terminal variance:
    \[
    \widehat{\Var}_g^1 = \int (z-\hat m_g^1)^2\,\hat p_1^g(\dd z);
    \]
    \item terminal mass:
    \[
    \hat M_g^1 = \mu_1^g(\Zcal);
    \]
    \item finite-measure discrepancy:
    \[
    D_g = \mathrm{Sinkhorn}^{\mathrm{unb}}(\mu_1^g,\nu_1^g).
    \]
\end{enumerate}

For the two-screen benchmark, these should be computed within each screen as well.

\subsection{Pass criteria for the core benchmark}

The model should satisfy:
\begin{enumerate}[label=(\roman*),leftmargin=2.2em]
    \item \(\hat m_{\mathrm{drift}}^1 > \hat m_{\mathrm{ctrl}}^1\);
    \item \(\widehat{\Var}_{\mathrm{diff}}^1 > \widehat{\Var}_{\mathrm{ctrl}}^1\);
    \item \(\hat M_{\mathrm{react}}^1 < \hat M_{\mathrm{ctrl}}^1\).
\end{enumerate}
Moreover, the model should achieve low terminal finite-measure discrepancy for all perturbations without needing to collapse every effect into one mechanism.

\subsection{Pass criteria for the mean-field benchmark}

Let screen \(2\) have the larger initial driver abundance. Then the fitted mean-field model should capture:
\begin{enumerate}[label=(\roman*),leftmargin=2.2em]
    \item a stronger rightward shift of the control endpoint in screen \(2\) than in screen \(1\);
    \item corresponding screen-specific shifts for other non-driver perturbations;
    \item improved joint fit relative to a context-free ablation.
\end{enumerate}

\section{Necessary ablations}

The benchmark is useful only if ablations fail in interpretable ways.

\subsection{Ablation 1: no growth field}

Set
\[
r_g \equiv 0.
\]
Then the reaction perturbation should fail to match the true terminal mass difference. If it does not fail, the benchmark is not actually testing the reaction layer.

\subsection{Ablation 2: shared diffusion}

Force \(\sigma_g\) to be identical across perturbations. Then the diffusion perturbation should fail to reproduce the true increase in terminal spread.

\subsection{Ablation 3: normalized-distribution loss only}

Replace the finite-measure endpoint loss by a loss on the normalized endpoint distributions only. Then the reaction perturbation should fail because terminal mass information is no longer supervised coherently.

\subsection{Ablation 4: no context layer}

In the two-screen mean-field benchmark, set
\[
c_{s,t}\equiv 0
\]
in the fitted model. Then the model should fail to explain the cross-screen endpoint differences induced by the driver perturbation.

\section{A recommended noise ladder}

The minimum useful simulation should first be run in a clean setting and only then be made harder. A good progression is:
\begin{enumerate}[label=(N\arabic*),leftmargin=2.6em]
    \item noiseless guide-abundance masses and exact latent coordinates;
    \item finite cell sampling for \(P4\) and \(P60\) while masses remain exact;
    \item noisy guide-abundance masses, for example
    \[
    \tilde M_g^1 = M_g^1 e^{\varepsilon_g},
    \qquad
    \varepsilon_g\sim\mathcal{N}(0,\sigma_M^2);
    \]
    \item multiple screens with different context trajectories.
\end{enumerate}

This ladder ensures that one first validates the mathematical structure before testing robustness.

\section{What the simulation proves, and what it does not}

\subsection{What it can prove}

A successful result on these two benchmarks would show that:
\begin{enumerate}[label=(\roman*),leftmargin=2.2em]
    \item the control-anchored weighted SDE can recover distinct perturbation effects on movement, spread, and mass from endpoint-only supervision;
    \item finite-measure endpoint fitting is genuinely useful when state sample size is separated from perturbation mass;
    \item the mean-field context layer becomes operationally useful when the data contain cross-screen variation in global environment.
\end{enumerate}

\subsection{What it cannot prove}

Such a simulation does \emph{not} prove that the biological system is truly one-dimensional, linear, or governed by an Ornstein--Uhlenbeck law. It also does not prove that the learned continuous-time dynamics is uniquely identifiable from two snapshots. The simulation is a \emph{minimal consistency check} for the architecture, not a full biological validation.

\section{Concise simulator specification}

The minimum useful simulation can therefore be summarized as follows.

\begin{quote}
First, build a one-dimensional four-perturbation benchmark with a shared Gaussian initial distribution, perturbation-specific Ornstein--Uhlenbeck drift targets, perturbation-specific diffusion amplitudes, and perturbation-specific reaction rates. Sample fixed numbers of cells at $P4$ and $P60$, record masses separately, and define finite-measure endpoints. This is the minimum benchmark for the control-anchored weighted SDE itself. Second, add one context-driving perturbation and create two synthetic screens that differ only in the initial abundance of that driver. Define a screen-level context as a soft occupancy functional of the joint population and let that context feed into the drift. This is the minimum benchmark that genuinely tests the value of a mean-field context layer.
\end{quote}

\section{Final summary}

The minimum useful simulation for this modeling approach should not attempt to emulate the full biological richness of a perturb-seq tumour ecosystem. Instead, it should be the smallest benchmark in which each claimed component of the model is necessary.

That leads to the following conclusion:

\begin{enumerate}[label=(\arabic*),leftmargin=2.2em]
    \item The \textbf{minimum useful core benchmark} is a one-dimensional four-perturbation endpoint-only simulation with separate drift, diffusion, and reaction perturbations and with finite-measure supervision.
    \item The \textbf{minimum useful mean-field benchmark} is a two-screen extension of the core benchmark with one driver perturbation whose abundance changes the global context trajectory.
\end{enumerate}

Anything simpler than the first benchmark fails to test the weighted control-anchored SDE properly. Anything simpler than the second benchmark fails to test the mean-field claim properly.

\end{document}